% ---
% id: receiver/rem:cdirection
% kind: rem
% title: "The direction of the unclamped $c$ at odd and even photon numbers"
% summary: >-
%   An admitted negative c lowers the analytic bound at odd photon numbers and
%   can raise it at even ones, by at most (K−1)ε₂², second order on the declared
%   domain. Neither comparison is an enclosure of the exact bound.
% prominence: standard
% status: accepted
% uses: [def:tolmap, def:bounddata, rem:tolerancedomain]
% mentions: [thm:exact-and-evaluated, def:enclosed-weight, rem:ia]
% assumes: []
% refs: []
% provenance: newly-derived
% lean:
% ---

\begin{remark}
\label{rem:cdirection}
Definition~\ref{def:tolmap}\,(iv) admits $c\in[-\varepsilon_2,0)$ and feeds
that value to \eqref{eq:beta} unchanged. Since
$\Bnd(n)=1+(K-1)\gamma c^{\,n}-\sum_k\beta_kd_k^{\,n}$ is increasing in
$c^{\,n}$, and since the exact $c$ satisfies $c_{\mathrm{ex}}\ge0$, the effect
splits by the parity of $n$, and both sides are stated.
\begin{itemize}[leftmargin=1.6em]
\item \emph{Odd $n$.} $c^{\,n}<0\le c_{\mathrm{ex}}^{\,n}$, so the evaluated
  $\Bnd(n)$ is \emph{below} the exact one. Conservative, and by a first-order
  amount: measured against the value at $c=0$ the deficit is
  $(K-1)\gamma|c|^{\,n}\le(K-1)\varepsilon_2$ at $n=1$, which is a lower bound
  on the deficit against $\Bnd(n)$ at $c_{\mathrm{ex}}$, that one being
  $(K-1)\gamma\big(c_{\mathrm{ex}}^{\,n}+|c|^{\,n}\big)$.
\item \emph{Even $n\ge2$.} $c^{\,n}=|c|^{\,n}>0$, so the evaluated $\Bnd(n)$
  can \emph{exceed} the exact one, by at most
  $(K-1)\gamma\big(|c|^{\,n}-c_{\mathrm{ex}}^{\,n}\big)
  \le(K-1)|c|^{\,n}\le(K-1)\varepsilon_2^{\,2}$, using $\gamma\le1$,
  $|c|\le\varepsilon_2$ and $\varepsilon_2\le1$. The bound
  $|c|\le\varepsilon_2$ is Definition~\ref{def:tolmap}\,(iv) and no other rule,
  which is why that item decides the admission of $c$ on \emph{every} route,
  item~(v)'s route-(ii) admission at $\varepsilon_1$ not being a bound on $|c|$
  at all. The last step is where the declared upper bound of
  Remark~\ref{rem:tolerancedomain} is used, since
  $\varepsilon_2^{\,n}\le\varepsilon_2^{\,2}$ for even $n\ge2$ holds exactly on
  $\varepsilon_2\le1$ and fails above it. The overshoot is second order in
  $\varepsilon_2$ on the declared domain, and this is the one direction in
  which clamping $c$ to $0$ would have been the safer choice.
\end{itemize}
The contract fixes the unclamped rule because a clamped and an unclamped
evaluation are both defensible and they differ, so leaving the choice to an
implementation would leave the emitted number undetermined by the contract.
Neither statement above is an enclosure: they compare two evaluations of
\eqref{eq:beta}, and Theorem~\ref{thm:exact-and-evaluated}\,(ii) still governs the distance
from either to the exact quantity. The case arises in the floating-point
evaluation only: the rigorous evaluation takes its lower end of $c$ in
$[0,1]$ (Remark~\ref{rem:ia}, Definition~\ref{def:enclosed-weight}), so no
negative $c$ enters a weight it forms.
\end{remark}
